\clearpage\thispagestyle{empty}
\pdfbookmark[0]{Quellen zum Bericht über Hardings Entdeckung vom Oktober 1804}{harding1804-sources}
\begin{center}\textbf{Quellen zum Bericht über Hardings Entdeckung vom Oktober 1804}\end{center}
\noindent Diese KI-Übertragung in heutiges Deutsch folgt dem vollständigen Auszug aus dem Bericht von Franz Xaver von Zach in \href{https://gdz.sub.uni-goettingen.de/id/PPN236007467}{Gauß’ \textit{Werken}, Band~VI} (Göttingen, 1874), S.~250--253. Von Zach berichtet über die Entdeckung von Karl Ludwig Harding und zitiert Carl Friedrich Gauß und Heinrich Wilhelm Matthias Olbers. Der ursprüngliche Bericht erschien als Beitrag XXXI im Oktoberheft 1804 der \textit{Monatlichen Correspondenz}, Band~X, S.~371--385. Die Überschrift des Auszugs stammt von \href{https://www.digitale-sammlungen.de/en/view/bsb10538602?page=385}{S.~371}. Die entsprechende Textpassage erstreckt sich von \href{https://www.digitale-sammlungen.de/en/view/bsb10538602?page=393}{S.~379} bis \href{https://www.digitale-sammlungen.de/en/view/bsb10538602?page=399}{S.~385}.
\par\smallskip\noindent Olbers verweist auf \href{https://archive.org/download/histoireclestef00lalagoog/page/n146.jpg}{S.~119} und \href{https://archive.org/download/histoireclestef00lalagoog/page/n158.jpg}{S.~131} der von Jérôme de Lalande herausgegebenen \href{https://archive.org/download/histoireclestef00lalagoog/page/n12.jpg}{\textit{Histoire céleste française}}, Band~I (Paris, Imprimerie de la République, Jahr IX [1801]). \href{https://catalogue.bnf.fr/ark:/12148/cb30722242z}{Katalog der Bibliothèque nationale de France}.
\par\smallskip\noindent Übertragung und Satz wurden mit KI erstellt. Die Digitalisate der Zeitschrift stammen von der Bayerischen Staatsbibliothek. Das von Google digitalisierte Exemplar des Sternkatalogs ist im Internet Archive zugänglich.
\clearpage
